% Options for packages loaded elsewhere
\PassOptionsToPackage{unicode}{hyperref}
\PassOptionsToPackage{hyphens}{url}
%
\documentclass[
  11pt,
]{article}
\usepackage{amsmath,amssymb}
\usepackage{iftex}
\ifPDFTeX
  \usepackage[T1]{fontenc}
  \usepackage[utf8]{inputenc}
  \usepackage{textcomp} % provide euro and other symbols
\else % if luatex or xetex
  \usepackage{unicode-math} % this also loads fontspec
  \defaultfontfeatures{Scale=MatchLowercase}
  \defaultfontfeatures[\rmfamily]{Ligatures=TeX,Scale=1}
\fi
\usepackage{lmodern}
\ifPDFTeX\else
  % xetex/luatex font selection
\fi
% Use upquote if available, for straight quotes in verbatim environments
\IfFileExists{upquote.sty}{\usepackage{upquote}}{}
\IfFileExists{microtype.sty}{% use microtype if available
  \usepackage[]{microtype}
  \UseMicrotypeSet[protrusion]{basicmath} % disable protrusion for tt fonts
}{}
\makeatletter
\@ifundefined{KOMAClassName}{% if non-KOMA class
  \IfFileExists{parskip.sty}{%
    \usepackage{parskip}
  }{% else
    \setlength{\parindent}{0pt}
    \setlength{\parskip}{6pt plus 2pt minus 1pt}}
}{% if KOMA class
  \KOMAoptions{parskip=half}}
\makeatother
\usepackage{xcolor}
\usepackage[margin=25mm]{geometry}
\setlength{\emergencystretch}{3em} % prevent overfull lines
\providecommand{\tightlist}{%
  \setlength{\itemsep}{0pt}\setlength{\parskip}{0pt}}
\setcounter{secnumdepth}{-\maxdimen} % remove section numbering
\usepackage{amsmath,amssymb,mathtools,xurl}
\emergencystretch=3em
\ifLuaTeX
  \usepackage{selnolig}  % disable illegal ligatures
\fi
\usepackage{bookmark}
\IfFileExists{xurl.sty}{\usepackage{xurl}}{} % add URL line breaks if available
\urlstyle{same}
\hypersetup{
  pdftitle={A facet-exchange counterexample to the fixed-mass regular-simplex conjecture},
  pdfauthor={mxym/math research project (AI-assisted)},
  hidelinks,
  pdfcreator={LaTeX via pandoc}}

\title{A facet-exchange counterexample to the fixed-mass regular-simplex
conjecture}
\author{mxym/math research project (AI-assisted)}
\date{8 October 2026}

\begin{document}
\maketitle

\subsection{Abstract}\label{abstract}

We give an explicit family of positive four-cell mass vectors for which
no translated regular-tetrahedral Gaussian partition maximizes the sum
of squared first moments. This disproves the arbitrary-prescribed-mass
statement of Conjecture 1.16 in Steven Heilman's \emph{Stable Gaussian
Minimal Bubbles}, arXiv:1901.03934v1 (2019). The centering in its
Problem 1.15 adds a constant under the prescribed-mass constraints. For
a tetrahedral apex displaced toward one vertex, its two kinds of
Gaussian facet measure are strictly unequal. The difference of the
adjacent cell moments therefore has a nonzero tangential component.
Exchanging two equal-measure small balls gives a genuine partition with
exactly the same masses and strictly larger objective. We prove that
every regular-tetrahedral partition with this mass vector is congruent
to the displaced candidate, and that a measurable optimizer exists, so
the refutation is not vacuous.

The equal-mass tetrahedral conjecture is not refuted or proved here.
Related unequal-mass positive-noise results already exist; they are
distinguished from the first-moment objective. No worldwide priority or
external human-review assertion is made.

\subsection{1. Exact conjecture and
theorem}\label{exact-conjecture-and-theorem}

Let \(\gamma_d\) denote the standard Gaussian probability measure on
\(\mathbb R^d\), with density
\(\varphi_d(x)=(2\pi)^{-d/2}\exp(-\|x\|^2/2)\). For a measurable
partition \(\mathcal C=(C_1,C_2,C_3,C_4)\), define

\[
 p_i=\gamma_d(C_i),\qquad m_i=\int_{C_i}x\,d\gamma_d(x),
 \qquad \mathcal E(\mathcal C)=\sum_{i=1}^4\|m_i\|^2.       \tag{1}
\]

Boundaries of the polyhedral cells below have Gaussian measure zero and
may be assigned to the smallest label. We use Lebesgue integrals and
Euclidean norms throughout.

Problem 1.15 in arXiv:1901.03934v1 fixes positive masses \(a_i\) summing
to one, defines a common vector \(w\) from a translated regular-simplex
partition, and writes \(w^{(i)}=w/a_i\). Its objective, up to the common
positive factor \(\sqrt{\pi/2}\), is

\[
 \sum_i\left\|\int_{C_i}(x-w/a_i)\,d\gamma_d(x)\right\|^2
 =\sum_i\|m_i-w\|^2
 =\mathcal E(\mathcal C)+4\|w\|^2.                       \tag{2}
\]

The last equality uses \(\sum_i m_i=\int x\,d\gamma_d(x)=0\). Thus the
prescribed-mass maximizers are exactly those of (1), regardless of the
fixed centering vector. Conjecture 1.16 states that for \(m-1\le d\) the
maximizing sets in Problem 1.15 are simplicial cones over a regular
simplex. Its \(m=4,d=3\), arbitrary-positive-mass case is enough to
refute the statement as written.

Use the regular-tetrahedral vectors

\[
 v_1=(1,1,1),\quad v_2=(1,-1,-1),\quad
 v_3=(-1,1,-1),\quad v_4=(-1,-1,1).                       \tag{3}
\]

They sum to zero, have norm \(\sqrt3\), and every pair is separated by
\(\ell=2\sqrt2\). For each \(t>0\), put \(w_t=t v_1\) and

\[
 C_i(t)=\{x:\langle v_i,x-w_t\rangle\ge
                   \langle v_j,x-w_t\rangle\ \forall j\},\qquad
 p_i(t)=\gamma_3(C_i(t)).                                \tag{4}
\]

The mass vector is specified by the fixed Gaussian integrals (4), not by
numerical optimization. In particular \(t=1\) gives a single fixed
counterexample instance with \(w_1=(1,1,1)\).

\textbf{Theorem 1 (counterexample for every \(t>0\)).} For each mass
vector \((p_1(t),p_2(t),p_3(t),p_4(t))\) in (4):

\begin{enumerate}
\def\labelenumi{\arabic{enumi}.}
\tightlist
\item
  All entries are positive, \(p_1(t)<1/4\), and
  \(p_2(t)=p_3(t)=p_4(t)=(1-p_1(t))/3\).
\item
  There is a measurable four-cell partition with precisely these masses
  and strictly larger \(\mathcal E\) than \(\mathcal C(t)\).
\item
  Every regular-tetrahedral conical partition with precisely these
  masses, allowing any rotation and apex, has the same objective as
  \(\mathcal C(t)\).
\item
  A maximizing measurable partition with these masses exists. None of
  the maximizing partitions can be regular-tetrahedral conical.
\end{enumerate}

Consequently the arbitrary-mass Conjecture 1.16 just quoted is false.
The same counterexample persists by cylindrical extension in every
ambient dimension \(d\ge3\).

\textbf{Corollary 1.1 (a full interval of mass vectors).} For every
\(0<p<1/4\), regular-tetrahedral cones fail to maximize (1) at masses
\((p,(1-p)/3,(1-p)/3,(1-p)/3)\). In particular, counterexamples occur
arbitrarily close to the equal-mass vector. Indeed (5) and its null
boundaries show that \(p_1(t)\) is continuous, strictly decreases from
\(1/4\) to zero as \(t\) goes from zero to infinity, and hence has
precisely that range. Continuity and the endpoint limit follow by
dominated convergence. Theorem 1 supplies the strict improvement at
every such \(t\).

\subsection{2. Exact geometry and the two facet
measures}\label{exact-geometry-and-the-two-facet-measures}

Write \(\phi(u)=(2\pi)^{-1/2}e^{-u^2/2}\) and let \(\Phi\) be its
distribution function. Symmetries permuting \(v_2,v_3,v_4\) and fixing
\(v_1\) are orthogonal and fix \(w_t\). They give the three equal masses
in Theorem 1. The first cell is

\[
 C_1(t)=\{(x,y,z):x+y\ge2t,\ x+z\ge2t,\ y+z\ge2t\}.      \tag{5}
\]

At \(t=0\) symmetry gives mass \(1/4\) to every cell. For \(t>0\), (5)
is strictly smaller up to a set of positive Gaussian measure. The set
difference contains an open ball, for example around a point with all
three coordinates strictly between \(0\) and \(t\). Every cell contains
a nonempty open set, so all four masses are positive.

For an interface between cells \(i,j\), its Gaussian facet measure is

\[
 A_{ij}(t)=\int_{\Sigma_{ij}(t)}\varphi_3(x)\,d\mathcal H^2(x),
 \quad \Sigma_{ij}(t)=\overline C_i(t)\cap\overline C_j(t). \tag{6}
\]

Triple intersections have surface measure zero. Symmetry leaves two
values: \(A=A_{12}=A_{13}=A_{14}\) and \(B=A_{23}=A_{24}=A_{34}\). Both
are strictly positive.

\textbf{Lemma 2 (strict facet imbalance).} For every \(t>0\),

\[
 A(t)=\phi(\sqrt2t)\,\mathbb P(X\ge t+|Y|/\sqrt2),\qquad
 B(t)=\phi(0)\,\mathbb P(Z\le\min(X/\sqrt2,2t-X/\sqrt2)),\tag{7}
\]

where the Gaussian coordinates in each probability are independent
standard normals. In particular \(0<A(t)<A(0)=B(0)<B(t)\).

\emph{Proof.} The interface \(\Sigma_{12}\) lies in \(y+z=2t\). Use its
orthonormal coordinates \(X=x\), \(Y=(y-z)/\sqrt2\), so
\(y=t+Y/\sqrt2,z=t-Y/\sqrt2\). The two remaining inequalities in (5)
become \(X\ge t+|Y|/\sqrt2\). Its normal coordinate is \(\sqrt2t\),
proving the first formula in (7).

The interface \(\Sigma_{23}\) lies in \(x=y\). With orthonormal
coordinates \(X=(x+y)/\sqrt2\) and \(Z=z\), the conditions that the
common score of \(v_2,v_3\) dominates \(v_1\) and \(v_4\) are
respectively \(Z\le2t-X/\sqrt2\) and \(Z\le X/\sqrt2\). Its normal
coordinate is zero, proving the second formula.

The first probability strictly decreases as \(t\) grows from zero, and
\(\phi(\sqrt2t)<\phi(0)\). In the second formula the domain at \(t=0\)
is \(Z\le-|X|/\sqrt2\), contained in the domain for \(t>0\). The
inclusion is strict on an open set: take \(0<X<t/\sqrt2\) and \(Z\)
close to zero. Gaussian densities are positive there. Finally
\(A(0)=B(0)\) by tetrahedral orthogonal symmetry. This proves every
strict comparison. \(\square\)

\subsection{3. First moments and a tangential
obstruction}\label{first-moments-and-a-tangential-obstruction}

\textbf{Lemma 3 (Gaussian flux formula).} For the cells in (4),

\[
 m_i=\frac1\ell\sum_{j\ne i} A_{ij}(v_i-v_j).              \tag{8}
\]

\emph{Proof.} The Gaussian identity \(\nabla\varphi_3=-x\varphi_3\) and
the divergence theorem give the negative outward normal integral on each
cell. On its \(ij\) facet the inward normal is \((v_i-v_j)/\ell\). To
justify the unbounded cells, first intersect them with a radius-\(R\)
ball. The artificial spherical boundary integral has magnitude at most a
constant times \(R^2e^{-R^2/2}\), tending to zero. The moment and facet
integrals converge absolutely. This gives (8). \(\square\)

Using (3) and the two facet values,

\[
 m_1=\frac{4A}{\ell}v_1,\qquad
 m_j=\frac{(A+3B)v_j+(B-A)v_1}{\ell}\quad(j=2,3,4).       \tag{9}
\]

Let \(e=(1,0,0)\) and \(n=(0,1,1)/\sqrt2\); the first is tangent to
\(\Sigma_{12}\) and the second is its unit normal into \(C_1\). In these
orthonormal directions the moment difference is

\[
 d=m_1-m_2=\sqrt2(A-B)e+(3A+B)n.                         \tag{10}
\]

Thus \(d\cdot e<0\) by Lemma 2. A volume-constrained first variation
would require the scalar score \(d\cdot x\) to be constant along this
facet. It is not constant. The next section does not rely on a formal
variational argument: it constructs an actual exact-mass improvement.

\subsection{4. An actual equal-measure ball
exchange}\label{an-actual-equal-measure-ball-exchange}

Fix any \(t>0\) and its positive constants \(A<B\). Define

\[
 \eta=\frac{B-A}{100(A+B)}\in(0,1/100),\quad
 q=(t+1,t,t)-\eta n,\qquad r=(t+2,t,t)+\eta n.             \tag{11}
\]

The open ball \(U\) centered at \(q\) of radius \(\eta/2\) lies strictly
inside \(C_2(t)\), and the open ball \(V\) centered at \(r\) of the same
radius lies strictly inside \(C_1(t)\). Indeed their centers have
distance \(\eta\) from the \(12\) plane on their respective sides.
Distances from either center to the other relevant separating planes
exceed \(1/2\), by the explicit inequalities used in Section 2 and
\(\eta<1/100\). Moving by \(\eta/2\) cannot cross any of those planes.
The two balls are disjoint.

Put \(\delta=\tfrac12\min(\gamma_3(U),\gamma_3(V))>0\). Gaussian measure
of a concentric ball is continuous and strictly increasing in its radius
from zero. There are therefore smaller concentric balls \(U'\subset U\)
and \(V'\subset V\) with

\[
 \gamma_3(U')=\gamma_3(V')=\delta.                        \tag{12}
\]

Transfer \(U'\) from cell 2 to cell 1 and \(V'\) from cell 1 to cell 2;
leave cells 3 and 4 unchanged. This defines a measurable partition
\(\mathcal C'\) with exactly the same four masses. If

\[
 \Delta=\int_{U'}x\,d\gamma_3(x)-\int_{V'}x\,d\gamma_3(x), \tag{13}
\]

its changed moments are \(m_1+\Delta,m_2-\Delta\). We verify strict
improvement rather than assume it from the tangent calculation.

By (10), \(d\cdot(q-r)=\sqrt2(B-A)-2\eta(3A+B)\) and
\(\|d\|\le\sqrt2|A-B|+3A+B<5(A+B)\). For every \(x\in U'\) and
\(z\in V'\), Cauchy--Schwarz consequently gives

\begin{equation*}
\begin{aligned}
 d\cdot(x-z)
 &\ge\sqrt2(B-A)-2\eta(3A+B)-\eta\|d\|\\
 &>\left(\sqrt2-\frac6{100}-\frac5{100}\right)(B-A)
   >(B-A)>0.
\end{aligned}                                             \tag{14}
\end{equation*}

Here \(\sqrt2>7/5\) is sufficient for the last inequality. Integrating
(14) over the product of the two restricted Gaussian measures and using
their common mass \(\delta\) yields \(d\cdot\Delta>\delta(B-A)\).
Finally the exact norm identity gives

\[
 \mathcal E(\mathcal C')-\mathcal E(\mathcal C(t))
 =2d\cdot\Delta+2\|\Delta\|^2
 >2\delta(B-A)>0.                                        \tag{15}
\]

This proves part 2 of Theorem 1 with genuine Gaussian volumes and strict
inequalities. No floating computation or solver output enters.

\subsection{5. All regular-tetrahedral candidates with these
masses}\label{all-regular-tetrahedral-candidates-with-these-masses}

It is not enough to improve one chosen apex. We must rule out a
different regular tetrahedron with the same mass vector.

Every regular-tetrahedral conical partition, after an orthogonal change
of coordinates and common harmless scaling of scores, can be written

\[
 D_i(\lambda)=\{x:\langle v_i,x\rangle-\lambda_i
                     \ge\langle v_j,x\rangle-\lambda_j\ \forall j\},
 \qquad\sum_i\lambda_i=0.                               \tag{16}
\]

Its apex is the unique \(w\) with \(\lambda_i=\langle v_i,w\rangle\).
The linear map \(x\mapsto(\langle v_i,x\rangle)_i\) is a bijection from
\(\mathbb R^3\) onto the hyperplane \(\sum_i s_i=0\). The score vector
has a strictly positive Gaussian density on that hyperplane and is
invariant under every label permutation.

\textbf{Lemma 4 (strict mass-price ordering).} In (16),

\[
 \gamma_3(D_i(\lambda))=\gamma_3(D_j(\lambda))
           \quad\Longleftrightarrow\quad\lambda_i=\lambda_j.
                                                                  \tag{17}
\]

If \(\lambda_i<\lambda_j\), the first mass is strictly larger.

\emph{Proof.} Swapping labels \(i,j\) preserves the score distribution.
Swap their prices to obtain \(\lambda'\). Then the mass of cell \(i\) at
\(\lambda'\) is the mass of cell \(j\) at \(\lambda\). If
\(\lambda_i<\lambda_j\), changing from \(\lambda'\) to \(\lambda\)
lowers the price of score \(i\) and raises that of score \(j\), with
other prices fixed. Every score vector won by \(i\) before is still won
by \(i\) after. The inclusion is strict on an open subset of the score
hyperplane: take \(s_i,s_j\) equal and sufficiently large, and all other
scores sufficiently small, subject to zero sum. Before the change \(j\)
wins, and after it \(i\) wins. Small perturbations retain both strict
wins. The positive density gives strict mass increase. Equality of
prices gives equality of masses by symmetry. \(\square\)

The three equal masses \(p_2,p_3,p_4\) therefore force
\(\lambda_2=\lambda_3=\lambda_4\). Since their sum is zero, their prices
have form \((3u,-u,-u,-u)\) and their apex is \(u v_1\). Also
\(p_1<p_2\) forces \(u>0\). By (5), \(p_1(u)\) is a strictly decreasing
function of \(u\): for \(u_2>u_1\), the cell inclusion is strict on an
open ball centered at \((b,b,b)\) with \(u_1<b<u_2\). Thus
\(p_1(u)=p_1(t)\) forces \(u=t\). Orthogonal changes preserve all
moments and norms. This proves part 3, including rotations and
permutations of the three equal-mass labels.

\subsection{6. Existence of a genuine maximizing
partition}\label{existence-of-a-genuine-maximizing-partition}

For completeness we give the existence argument, so part 4 and the
refutation do not rely on an unproved assumption about attaining the
supremum. Fix any positive masses \(p_i\) summing to one.

Consider the fractional partitions in the real Hilbert space
\(L^2(\gamma_3)^4\):

\[
 \mathcal K=\{(f_i):0\le f_i\le1,\ \sum_i f_i=1
                \text{ a.e.},\ \int f_i\,d\gamma_3=p_i\}. \tag{18}
\]

This is a nonempty bounded norm-closed convex set, and is weakly compact
by Hilbert-space reflexivity. Each coordinate of
\(m_i(f)=\int x f_i(x)\,d\gamma_3(x)\) is a continuous linear functional
because each Gaussian coordinate belongs to \(L^2\). Thus the objective,
which depends on finitely many such coordinates, is weakly continuous
and has a maximum \(M\) at some \(f^*\). Let \(m_i^*=m_i(f^*)\).

The linear functional \(L(f)=\sum_i m_i^*\cdot m_i(f)\) satisfies
\(L(f)\le\sqrt M\sqrt{\mathcal E(f)}\le M\), with equality at \(f^*\).
Its maximizers form a nonempty weakly compact face of \(\mathcal K\).
The Krein--Milman theorem supplies an extreme point of this face, which
is also extreme in \(\mathcal K\).

Every extreme point of \(\mathcal K\) is an indicator partition.
Otherwise two components lie between \(\epsilon\) and \(1-\epsilon\) on
a common positive-measure region, for some \(\epsilon>0\) (use the
finite number of component pairs and a countable union over positive
rational \(\epsilon\)). Gaussian measure is nonatomic. Split that region
into two disjoint sets of equal positive measure and put
\(h=1_{S_+}-1_{S_-}\). Adding \(\epsilon h/2\) to one component and
subtracting it from the other, or reversing both changes, yields two
distinct elements of \(\mathcal K\) with the original as midpoint. This
contradicts extremality.

Call the indicator extreme point \(g\). Its \(L(g)=M\) and
Cauchy--Schwarz imply \(\mathcal E(g)\ge M\) when \(M>0\). If \(M=0\),
every member already has the maximum; a measurable partition with the
prescribed masses is obtained by splitting an atomless probability
space. In either case a genuine maximizing partition exists. The
standard results used here are Hilbert-space reflexivity, weak
compactness of bounded closed convex subsets of a reflexive space, and
the Krein--Milman theorem; none asserts Gaussian simplex optimality.

Part 2 gives a strictly better partition than the objective of every
regular-tetrahedral candidate identified in part 3. The existing
optimizer cannot be any of those candidates. This finishes Theorem 1.
For \(d>3\), append independent unused Gaussian coordinates to the
construction. Orthogonal reduction of any regular four-cell simplex to
its three-dimensional span gives the same uniqueness argument.
\(\square\)

\subsection{7. Literature and exact scope of the
refutation}\label{literature-and-exact-scope-of-the-refutation}

The target is the \textbf{arbitrary-positive-prescribed-mass statement
as written} in Heilman, \emph{Stable Gaussian Minimal Bubbles},
arXiv:1901.03934v1, Section 1.4, Problem 1.15 and Conjecture 1.16,
printed page 8 (13 January 2019). Its equal-mass specialization is
different and is untouched by the present counterexample.

Heilman--Mossel--Neeman, \emph{Standard Simplices and Pluralities are
Not the Most Noise Stable}, Israel Journal of Mathematics 213 (2016),
33--53, arXiv:1403.0885v3, Theorem 2.6, already proves a nonoptimality
statement for unequal masses and every positive noise parameter
\(0<\rho<1\). That theorem statement does not itself assert strict
improvement of the first-Hermite coefficient: a strictly positive
function for \(\rho>0\) may have zero derivative at zero. Here strict
first-moment improvement is proved directly by (15). No novelty claim is
based merely on taking a limit of the known noisy result.

The finite literature comparison is in \texttt{PRIOR\_WORK.md}; the
archived primary-source hashes are in \texttt{SOURCES.json}. We do not
claim this obstruction has never been observed. The proof uses Gaussian
integrals and elementary inequalities, with no computational theorem
input. The partial Lean check verifies algebraic identities and the
strict exchange-margin reduction, not the full Gaussian geometry or
existence theorem. All mathematical claims above have explicit written
proofs.

\end{document}
